\pagebreak
\setcounter{section}{10}
\chapter{Formation and Evolution of LMXBs}
\hrule
\vspace{.5cm}

\section{Rare formation process}
\begin{list}{-}{}
\item Due to the \textit{lower} mass of the donor star when compared to the accretor, mass transfer in LMXBs can be long stable on large timescales, such as \dots
\begin{list}{-}{}
\item nuclear (fuel of the star)
\item orbital shrinking due to loss of angular momentum
\end{list}
\item When comapring the theoretical lifespan of them, to the quanity obsreved, we find that they have formation rate of \( 10^{-7} \text{yr}^{-1}\)
\item one out of \(\sim 10^{5}\) finishes in a LMXB system, whereas HMXBS it is \textit{one out of ten} 
\end{list}

\section{Possible formation channels}
\begin{list}{-}{}
\item \gls{ac:AIC} of a massive WD reaching \gls{chandrasekhar-limit} \ref{sec:LMXBAICFormation}
\item Formation Via Common Envelope Evolution \ref{sec:LMXBFormViaCE}
\end{list}
\pagebreak
\subsection{Formation Via Common Envelope Evolution\label{sec:LMXBFormViaCE}}
\hrule
\vspace{.5cm}
Deduced by looking at the Her x-1 system 
\begin{list}{-}{}
\item far distance (3kpc) from the galactic plane, but a family young star
\item must have experienced something to launch hit out of the plane at 120kms\(^{-1}\)
\item This was from an SN, most likely some sort of \gls{natal-kick}
\item The velocity of this NS is much greater than a possible value formed via \gls{ac:AIC} on a WD
\end{list}
\textbf{Evolutionary Steps}
In this config the initial system consists of a relatively massive star (\(\sim 12 \text{to} 15 M_\odot\)) together with a \(2M_\odot\) star in a wide orbit (\(P_{orb} > 1yr\)) 

\subsection{\gls{ac:AIC} of a massive WD reaching \gls{chandrasekhar-limit} \label{sec:LMXBAICFormation}}

\subsection{Origin from a HMXB with a much lower mass companion (triple star system)}




